% Generated draft: requires mathematical and source review.
\documentclass[11pt]{article}
\usepackage[margin=25mm]{geometry}
\usepackage{fontspec}
\setmainfont{Latin Modern Roman}
\usepackage{amsmath,amssymb,mathtools,unicode-math}
\setmathfont{Latin Modern Math}
\usepackage{enumitem,needspace,xurl,hyperref}
\setlength{\parindent}{0pt}
\setlength{\parskip}{5pt}
\setlength{\emergencystretch}{3em}
\newcommand{\N}{\mathbb{N}}
\newcommand{\Z}{\mathbb{Z}}
\newcommand{\Q}{\mathbb{Q}}
\newcommand{\R}{\mathbb{R}}
\newcommand{\C}{\mathbb{C}}
\newcommand{\F}{\mathbb{F}}
\newcommand{\A}{\mathbb{A}}
\newcommand{\PP}{\mathbb{P}}
\newcommand{\E}{\mathbb{E}}
\newcommand{\eps}{\varepsilon}
\newcommand{\dd}{\,\mathrm{d}}
\begin{document}
\section*{UnsolvedMath ID: 800516}
\textbf{Source problem code:} catalogue entry
\subsection{Definitions and mathematical statement}
Let \(n,k\in\N\) satisfy \(1<k<n\), and let \(A\in\C^{n\times n}\) be non-Hermitian, meaning \(A\ne A^*\), where \(^*\) denotes conjugate transpose. For \(Q\in\C^{n\times k}\) with \(Q^*Q=I_k\), define the compression \(B_Q=Q^*AQ\). Its eigenvalues \(\theta_1(Q),\ldots,\theta_k(Q)\), counted with algebraic multiplicity, are the Ritz values associated with \(Q\). Define the numerical range by
\[
W(A)=\{v^*Av:v\in\C^n,\ \|v\|_2=1\}.
\]
The supplied source recalls that \(W(A)\) is closed and convex, contains the eigenvalues of \(A\), and contains every Ritz value; it may also contain points far from the eigenvalues of \(A\). To express the distribution of the Ritz values, one may use the empirical measure
\[
\mu_{A,Q}=\frac1k\sum_{j=1}^k\delta_{\theta_j(Q)},
\]
where \(\delta_z\) is the unit point mass at \(z\).

\textbf{Source question.} What deterministic or probabilistic information can be obtained about the distribution of these Ritz values across \(W(A)\) when \(1<k<n\)? How does the question change when the columns of \(Q\) form an orthonormal basis of a Krylov subspace, or an orthonormal basis of a random subspace?

For clarity, a standard single-vector Krylov subspace is
\[
\mathcal K_k(A,b)=\operatorname{span}\{b,Ab,\ldots,A^{k-1}b\},\qquad b\in\C^n\setminus\{0\}.
\]
Under this interpretation, an orthonormal basis with exactly \(k\) columns requires \(\dim\mathcal K_k(A,b)=k\). The source does not specify the starting vector, whether this particular Krylov construction is intended, or how breakdown is to be handled. It also does not specify a probability law for a random subspace, any randomness in the Krylov construction, a particular distributional statistic or bound, or an asymptotic regime. Thus the supplied formulation is an open research question rather than a fully specified conjecture; its probabilistic variants require additional choices.
\subsection{Short English statement}
For a non-Hermitian matrix, describe how the eigenvalues of its compressions to subspaces of dimension \(k\), with \(1<k<n\), are distributed inside its numerical range. Investigate arbitrary subspaces, Krylov subspaces, and random subspaces. The source leaves the desired distributional conclusions, Krylov construction, and probability laws unspecified.
\subsection{Sources}
\begin{itemize}
\item Local UnsolvedMath registry, supplied complete source record: permanent ID \(800516\), source problem code \texttt{catalogue entry}, title ``Linear Systems and Eigenvalue Problems --- Problem 25.'' The registry attributes the problem to Amsel, Baumann et al., with proposed year \(2026\), and labels it \texttt{open}. This status is supplied metadata, not an independently verified mathematical or literature assessment.
\item Amsel, Baumann et al., \emph{Linear Systems and Eigenvalue Problems: Open Questions from a Simons Workshop} (2026), Problem 25. Supplied source URL: \url{https://arxiv.org/abs/2602.05394}. The provenance records access on \texttt{2026-10-16} and extraction by \texttt{arxiv\_latex}; these are supplied provenance statements, not a claim of independent source access.
\item Supplied upstream revision: \texttt{d371ee917ac9647fbb06cd6fc529deec6f3dda82}.
\end{itemize}
\end{document}
